\documentclass[12pt]{article}
\usepackage[utf8]{inputenc}
\usepackage[T1]{fontenc}
\usepackage{hyperref}
\title{Tacana Language Resources for K-12 Curriculum Development}
\author{Compiled by Claude Sonnet 5 (claude-sonnet-5) for Andr\'as Kornai}
\date{}
\begin{document}
\maketitle

\section{Introduction}
Tacana (ISO 639-3 tna) is an endangered Tacanan language of Bolivia, with about 1,800 speakers. This report rests on a single round of web searches and found only a few resources.

\section{Primary Educational Resources}
No primers, textbooks or school programmes were found.

\section{Linguistic and Reference Materials}
\begin{itemize}
\item Living Dictionaries Tacana dictionary. It is under construction, and community members are invited to correct spellings and record audio: \url{https://livingdictionaries.app/tacana}, \url{https://livingdictionaries.app/tacana/about}
\item SAPHON phonological inventory (UC Berkeley): \url{https://linguistics.berkeley.edu/saphon/en/inv/Tacana.html}
\item A bibliography cites \emph{Diccionario tacana-castellano, castellano-tacana} (1989) and Ottaviano's structural grammar of Tacana (1967, SIL). I saw only citations, not the works themselves, so no URL is given.
\end{itemize}

\section{Religious and Community Resources}
None found.

\section{Digital and Audio Resources}
The Living Dictionaries project is the only digital resource found. It may include audio, which I did not verify.

\section{Curriculum Development Context}
Tacana has a dictionary in progress and older SIL-era descriptions, but no pedagogical materials were found. A curriculum would need new primers built on the existing dictionary and orthography work, in partnership with the community. Araona and Cavine\~na are related languages, but they are different languages and their materials should not be substituted.
\end{document}